\uuid{bG74}
\titre{Fraction}
\auteur{bodin ; hanani ; mzari}
\organisation{exo7}
\licence{CC BY-NC-SA 4.0}
\datecreate{2020-01-20}
\niveau{L1}
\module{Analyse}
\chapitre{Propriétés de R}
\sousChapitre{Les rationnels}
\difficulte{2}
\theme{Rationnels, fraction}
\exo7qcmid{lille-1.reels.004}

\contenu{
\begin{question}
\qclassification{120.01}
\qtags{motcle=fraction}
Quelles sont les assertions vraies ?
\begin{answers}
    \good{La somme de deux nombres rationnels est un nombre rationnel.}

    \good{Le produit de deux nombres rationnels est un nombre rationnel.}

    \bad{La somme de deux nombres irrationnels est un nombre irrationnel.}

    \bad{Le produit de deux nombres irrationnels est un nombre irrationnel.}
\end{answers}
\begin{explanations}
La somme de deux nombres rationnels est un nombre rationnel. Le produit aussi.
C'est en général faux pour les nombres irrationnels !
\end{explanations}
\end{question}
}
